% SGA 3, 第 VI_B 篇, 第 1 轮英文重构.
% 范围: 推论 4.3 与 4.4 全文, 含编者注 31 和 32.
% 底本: Polo--Gille Exp6B-13oct24.pdf, 分篇第 24--25 页
% (合订本第 365--366 页; 原版页边印刷页码 350--351),
% 止于第 5 节标题之前.
% OCR 与 Reinhold 英文 Markdown 仅为比对材料.

\medskip
\phantomsection\label{cor:sga3-VIB-4.3}%
\noindent\textbf{推论 4.3.---}
\emph{设 \(G\) 是局部有限呈示的平坦 \(S\)-群. 则函数}
\[
  s\longmapsto\dim G_s
\]
\emph{在 \(S\) 上局部常值.}

这立即由 \SGARefLink{prop:sga3-VIB-4.2}{4.2} 得到,
因为每个局部有限呈示的平坦态射都普遍开 (EGA~IV\(_2\), 2.4.6).

\medskip
\phantomsection\label{cor:sga3-VIB-4.4}%
\noindent\textbf{推论 4.4.---}
\emph{设 \(G\) 是在单位截面的各点处对 \(S\) 局部有限呈示的
\(S\)-群. 考虑下列条件:}
\begin{enumerate}
\item[\textup{\SGAInlineRefTarget{cor:sga3-VIB-4.4:item:i}(i)}]
  \emph{\(G\) 在单位截面的各点处对 \(S\) 平滑
  \textup{(参见 \hyperref[thm:sga3-VIB-3.10]{3.10})}.}
\item[\textup{\SGAInlineRefTarget{cor:sga3-VIB-4.4:item:ii}(ii)}]\phantomsection\label{cond:sga3-VIB-4.4-ii}
  \emph{对每个 \(s\in S\), \(G_s\) 在 \(\kappa(s)\) 上平滑, 且函数}
  \[
    s\longmapsto\dim G_s
  \]
  \emph{在 \(S\) 上局部常值.}
\item[\textup{(iii)}] \emph{对每个 \(s\in S\),
  \(G_s\) 在 \(\kappa(s)\) 上平滑, 且单位截面有一个邻域 \(V\), 使得}
  \[
    \pi_V:V\longrightarrow S
  \]
  \emph{普遍开.}
\item[\textup{\SGAInlineRefTarget{cor:sga3-VIB-4.4:item:iv}(iv)}]
  \emph{对每个 \(s\in S\), \(G_s^0\) 在 \(\kappa(s)\) 上平滑,
  且 \(G^0\) 可由 \(G\) 的一个在 \(S\) 上普遍开的开子群概形表示.}
\end{enumerate}
\emph{则有下列蕴含关系:}
\[
  \textup{(i)}\Longrightarrow\textup{(ii)}
  \Longleftrightarrow\textup{(iii)}
  \Longleftrightarrow\textup{(iv)}.
\]
\emph{若 \(S\) 还约化, 则条件
\textup{\SGARefLink{cor:sga3-VIB-4.4:item:i}{(i)}}--\textup{\SGARefLink{cor:sga3-VIB-4.4:item:iv}{(iv)}}
等价, 并蕴含 \(G^0\) 在 \(S\) 上平滑.}%
\footnote{\label{note:sga3-VIB-31}编者注: 这里不能省去 \(S\) 约化的假设.
事实上, 设 \(k\) 是域, \(S=\operatorname{Spec}A\),
其中 \(A=k[\delta]\), \(\delta^2=0\), 并设 \(G\) 是
\(\mathbb G_{a,S}\) 的闭子群
\[
  \operatorname{Spec}A[X]/(\delta X)
\]
(它对每个 \(A\)-代数 \(R\) 配给由满足 \(\delta r=0\) 的
\(r\in R\) 所成的子群). 则 \(G\) 是满足
\textup{\SGARefLink{cor:sga3-VIB-4.4:item:ii}{(ii)}}--\textup{\SGARefLink{cor:sga3-VIB-4.4:item:iv}{(iv)}}
的 \(S\)-群, 但它在 \(S\) 上不平坦, 因而不平滑.}

证明 \textup{\SGARefLink{cor:sga3-VIB-4.4:item:i}{(i)}} 蕴含
\textup{\SGARefLink{cor:sga3-VIB-4.4:item:ii}{(ii)}}.
对每个 \(x\in G\), 由 \hyperref[lem:sga3-VIB-1.5]{1.5},
\[
  \dim_x\pi^{-1}\bigl(\pi(x)\bigr)
    =\dim\pi^{-1}\bigl(\pi(x)\bigr).
\]
因此由 EGA~IV\(_4\), 17.10.2, 函数
\[
  x\longmapsto
  \dim_x\pi^{-1}\bigl(\pi(x)\bigr)
    =\dim\pi^{-1}\bigl(\pi(x)\bigr)
\]
在单位截面的某邻域上连续.
故函数 \(s\mapsto\dim G_s\) 在 \(S\) 上连续,
从而在 \(S\) 上局部常值.
另一方面, 由 \SGARefLink{cor:sga3-VIB-1.3.1}{1.3.1},
对每个 \(s\in S\), \(G_s\) 在 \(\kappa(s)\) 上平滑.

证明 \textup{\SGARefLink{cor:sga3-VIB-1.3.1:item:ii}{(ii)}} 蕴含
\textup{\SGARefLink{cor:sga3-VIB-4.4:item:iv}{(iv)}}.
只须证明 \(\underline{G^0}\) 在 \(G\) 中开;
因为此时由 \hyperref[case:sga3-VIB-3.4]{3.4},
\(G^0\) 可由 \(\underline{G^0}\) 上诱导的子群概形表示,
而上面所述 \(G^0\) 的性质将由 \hyperref[cor:sga3-VIB-2.4]{2.4}
和 \hyperref[prop:sga3-VIB-4.2]{4.2} 得到.
给定 \(s\in S\), 如 \hyperref[thm:sga3-VIB-3.10]{3.10} 的证明中那样构造
\[
  W,\quad T,\quad \pi'',\quad \epsilon'',\quad\text{和}\quad W^0.
\]
则由 EGA~IV\(_3\), 15.6.7, \(W^0\) 在 \(W\) 中开.
另一方面, 在假设 \hyperref[cond:sga3-VIB-4.4-ii]{4.4}\textup{(ii)} 下,
由 \hyperref[prop:sga3-VIB-4.2]{4.2} 可知 \(\pi\) 在 \(W^0\) 的每个点处普遍开.
因此由 VI\(_A\), \SGARefLink{sga3:VIA:0.1}{0.1},
\(\mu\) 在
\[
  W^0\times_S W^0,
\]
的每个点处普遍开, 这表明 \(W^0\cdot W^0\) 开;
如\hyperref[thm:sga3-VIB-3.10]{定理~3.10} 的证明那样完成证明.

显然
\[
  \textup{(iv)}\Longrightarrow\textup{(iii)},
\]
而 \(\textup{(iii)}\Longrightarrow\textup{(ii)}\)
由对 \(V\) 应用 \SGARefLink{prop:sga3-VIB-4.2}{4.2} 得到.

最后, 假设 \textup{\SGARefLink{prop:sga3-VIB-4.2:item:ii}{(ii)}}--\textup{\SGARefLink{thm:sga3-VIB-3.10:item:iv}{(iv)}}
成立, 来证明当 \(S\) 约化时, \(G^0\) 在 \(S\) 上平滑.%
\footnote{\label{note:sga3-VIB-32}编者注: 下面的论证扩充了原文.}
为此可假设 \(G=G^0\). 则由下面的 \SGARefLink{cor:sga3-VIB-5.5}{5.5},
\(G\) 在 \(S\) 上有限呈示. 因此 \(\pi_G\) 有限呈示,
其纤维几何整, 且相对维数在 \(S\) 上局部常值.
由 EGA~IV\(_3\), 15.6.7 可知, \(\pi_G\) 诱导的态射
\[
  G\times_S S_{\mathrm{red}}\longrightarrow S_{\mathrm{red}}
\]
平坦. 所以当 \(S\) 约化时, \(\pi_G\) 平坦.
在此情形下, 由\hyperref[thm:sga3-VIB-3.10]{定理~3.10},
\(G=G^0\) 在 \(S\) 上平滑.

% 翻译覆盖：本组件全文；源文件 39_expose_VIB_cor43_cor44_en.tex。
% 自检：行内数学 78，陈列数学 11，原生 tikzcd 0；全部数学及编号保留，公式内自然语言已译。
% 引用键核对：label 5，hyperref 10，SGARefLink 15，行内目标 3；脚注 2，星号脚注正文 0，均与源文件一致。
% 编译及渲染审查由主线程整合完成；此处未宣称编译或独立校对完成。
